\subsection{Afbeeldingen en hun eigenschappen}

\begin{lemma}[3.1.3]
    \begin{enumerate}
        \item[(i)] Een afbeelding $f \colon A \to B$ is bijectief \(\iff\) ze injectief én surjectief is.
        \item[(ii)] Als $f$ een bijectie is van $A$ naar $B$, dan is $|A| = |B|$, d.w.z.\ $A$ en $B$ hebben evenveel elementen.
    \end{enumerate}
\end{lemma}
\begin{proof}(i)
    Volgt uit definities 3.1.1.
\end{proof}
\begin{proof}(ii)
    Volgt uit definitie 3.1.1 (vi).
\end{proof}
\subsection{Kern en beeld van een lineaire afbeelding}

\begin{lemma}[3.2.2]
    Zij $f \colon V \to W$ een lineaire afbeelding. Dan is $\ker f \subseteq V$ een deelruimte van $V$ en $\operatorname{im} f \subseteq W$ een deelruimte van $W$.
\end{lemma}
\begin{proof}
    \(\ker (f) = \{ v \in V |f(v) = 0\}. \)
    \(\ker (f)\) is niet ledig want \(f(0) = 0\).
    
    Als \(v_1\) en \(v_2 \in \ker (f):\)
    \(f(\lambda v_1 + \mu v_2) = \lambda f( v_1) + \mu f( v_2)
    = \lambda 0 + \mu 0
    = 0\).

    \(\implies \lambda v_1 + \mu v_2 \in \ker (f)\).

    Als \(w_3\) en \(w_4 \in \operatorname{im} f: f(v_3) = w_3\) en \(f(v_4) = w_4 \in W\) 
    waarbij \(v_3,v_4 \in V\).

    Ook \(\operatorname{im}(f)\) is niet leeg want: \(f(0_V) = 0_W\).

    Dan 
    \(\lambda w_3 + \mu w_4 
    = \lambda f( v_3) + \mu f( v_4)
    =f(\lambda v_3 + \mu v_4) \)
    \(\implies \lambda w_3 + \mu w_4  \in \operatorname{im}(f) \).
\end{proof}

\begin{lemma}[3.2.4]
    Zij $f \colon V \to W$ een lineaire afbeelding. Dan geldt:
    \begin{enumerate}
        \item[(i)] $f$ is injectief als en slechts als $\ker f = \{0\}$.
        \item[(ii)] $f$ is surjectief als en slechts als $\operatorname{im} f = W$.
    \end{enumerate}
\end{lemma}
\begin{proof}(i)
    \([\Rightarrow]\) Stel \(f = \) injectief:

    Als \(f(v) = 0\), maar \(f(0) = 0\)
    \(\implies v = 0\).

    \([\Leftarrow]\) Stel omgekeerd \(\ker (f) = \{0\}\):

    Stel 
    \(
        f(v) = f(w)
         \implies f(v) - f(w) = 0 
            \implies f(v-w) = 0
         \implies v-w = 0
         \implies v = w.
    \)

    Dit is per definitie injectief.
\end{proof}
\begin{proof}(ii)
    \([\Rightarrow]\) Stel \(f = \) surjectief:

    Dit betekend: \(\forall w \in W \, \exists v \in V: f(v) = w\).
    Dus \(\{w \in W |\, \exists v \in V: f(v) = w \in W \} = W\).

    \([\Leftarrow]\) Stel omgekeerd \(\operatorname{im} (f) = W\):

    Per definitie van \(\operatorname{im}(f),\, \forall w \in W \,\exists v \in V: f(v) = w \).
    Dit is per definitie surjectief.
\end{proof}

\begin{lemma}[3.2.5]
    Zij $f \colon V \to W$ een lineaire afbeelding en zij $U \subseteq V$ een deelruimte. Beschouw de restrictie $f|_U$ van $f$ tot $U$. Dan is
    \[ \ker f|_U = \ker f \cap U \quad \text{en} \quad \operatorname{im} f|_U \subseteq \operatorname{im} f. \]
\end{lemma}
\begin{proof}
    \(\ker f|_U  = \{u \in U |f(u)=0\} = \{v \in V |f(v)=0\} \cap U= \ker f \cap U\).


    \(\operatorname{im}f|_U = \{f(u)|u\in U\} \subseteq \{f(v)|v\in V\} = \operatorname{im} f \).
\end{proof}

\begin{lemma}[3.2.6]
    Zij $V$ en $W$ twee $K$-vectorruimten, en zij $\mathcal{B}$ een basis voor $V$. Zij $f \colon V \to W$ een lineaire afbeelding. Dan geldt:
    \begin{enumerate}
        \item[(i)] $f$ is volledig bepaald door de beelden van de basiselementen $f(b)$ voor alle $b \in \mathcal{B}$;
        \item[(ii)] $\operatorname{im}(f) = \operatorname{span}(\{f(b) \mid b \in \mathcal{B}\})$.
    \end{enumerate}
\end{lemma}
\begin{proof}(i)
    \(v = \sum_i \mu_i b_i\) is uniek bepaald omdat \(\mathcal{B}\) een basis is.
    \(
        f(v) = f\left(\sum_i \mu_i b_i\right) = \sum_i \mu_i f(b_i).
    \)

    Wegens de uniciteit - \(f(v)\) ligt uniek vast \(\forall v \in V\) - volgt het te bewijzen.
\end{proof}
\begin{proof}(ii)
    \(
        \operatorname{im}(f)
        := \{f(v)|v\in V\}
        = \{\sum_i \mu_i f(b_i)| \mu_i \in K\}
        := \operatorname{span}(\{f(b) \mid b \in \mathcal{B}\}).
    \)
\end{proof}
\subsection{Injectiviteit en surjectiviteit}

\begin{stelling}[3.3.1]
    Zij $V, W$ twee $K$-vectorruimten, zij $\mathcal{B}$ een basis voor $V$, en zij $f \colon V \to W$ een lineaire afbeelding. Dan geldt:
    \begin{enumerate}
        \item[(i)] $f$ is injectief als en slechts als $\{f(b) \mid b \in \mathcal{B}\}$ een lineair onafhankelijke verzameling is in $W$ waarbij de elementen $f(b)$ twee aan twee verschillend zijn van elkaar;
        \item[(ii)] $f$ is surjectief als en slechts als $\{f(b) \mid b \in \mathcal{B}\}$ een voortbrengende verzameling is in $W$;
        \item[(iii)] $f$ is bijectief als en slechts als $\{f(b) \mid b \in \mathcal{B}\}$ een basis is in $W$ waarbij de elementen $f(b)$ twee aan twee verschillend zijn van elkaar.
    \end{enumerate}
\end{stelling}
\begin{proof}(i)
    \([\Rightarrow]\) Stel \(f\) is injectief: dus \(\ker f = \{0\}\).

     \(f(v) = \sum_i \mu_i f(b_i) = f(\sum_i \mu_i b_i) =0 \iff \sum_i \mu_i b_i = 0.\)

     Omdat \(\mathcal{B}\) een basis is, is \(\forall \mu \in K, \, \mu_i = 0\). Dus lineair onafhankelijk.

    \([\Leftarrow]\) Stel omgekeerd \(\{f(b) \mid b \in \mathcal{B}\}\) lineair onafhankelijk:

    Als \(f(v_1) = f(v_2)\):
    \[
        \sum_{i} \mu_i f(b_i) = \sum_{i} \nu_i f(b_i)
        \implies 0= \sum_{i} \mu_i f(b_i) - \sum_{i} \nu_i f(b_i)
        = \sum_{i} \left( \mu_i  -  \nu_i \right) f(b_i).
    \]
    
    Wegens lineair onafhankelijk zijn is \(\mu_i = \nu_i\).
    \(  f(v_1) = f(v_2)  \implies v_1 = v_2\).
\end{proof}
\begin{proof}(ii)

     \([\Rightarrow]\) Stel \(f\) surjectief:

     \(\forall w \in W \, \exists v \in V: f(v) = w 
     \implies \forall w \in W: w = f(v) = \sum_{i} \mu_i f(b_i)\).
     \begin{align*}
        \operatorname{span}(f(b)) 
        &= \{\sum_{i} \mu_i f(b_i) | \mu_i \in K\}\\
        &= \{w | w \in W\} = W\\
     \end{align*}
    \([\Leftarrow]\) Stel omgekeerd \(\{f(b) \mid b \in \mathcal{B}\}\) voortbrengend voor \(W\):
    \(
    \operatorname{span}(\{f(b)\mid b \in \mathcal{B}\}) = W.
    \)

    Volgens lemma 3.2.6 (ii) is \(\operatorname{span}(\{f(b)\mid b \in \mathcal{B}\}) = \operatorname{im}(f)\implies \operatorname{im}(f) =W\).
    
    \(\operatorname{im}f := \{w \in W | \exists v \in V: f(v) = w\} = W\).
     Wat de definitie van surjectiviteit is.
\end{proof}
\begin{proof}(iii)
    Volgt uit (i) en (ii).
\end{proof}

\begin{gevolg}[3.3.5]
    Twee $K$-vectorruimten zijn isomorf als en slechts als ze dezelfde dimensie hebben.
\end{gevolg}
\begin{proof}
    Isomorfisme := bijectieve lineaire afbeelding.

    \([\Rightarrow]\) Stel \(W \cong V\):

    Wegens stelling 3.3.1 is er een basis voor \(V\) die ook een basis is voor \(W\).
     \(\implies \dim(V) = |\mathcal{B}| = \dim(W)\).

    \([\Leftarrow]\) Omgekeerd stel \(\dim(V) = \dim(W)\):

    \(\dim(V) = \dim(W) = |\mathcal{B}|=|\mathcal{C}|\)
    met \(|\mathcal{B}|\) een basis voor \(V\) en \(|\mathcal{C}|\) een basis voor \(W\).

    We definieren een lin afb: \(f \colon V \to W: \sum_i \mu_i b_i \to \sum_i \mu_i \beta(b_i)\)

    Wegens ze evenveel elementen bevatten zijn ze 2 aan 2 verschillend. 
    Uit stelling 3.3.1 volgt nu dat \(f\) een bijectie dus isomorfisme is.
\end{proof}

\begin{gevolg}[3.3.6]
    Zij $K$ een veld, en zij $V$ een eindig-dimensionale $K$-vectorruimte. Stel $n = \dim_K(V)$. Dan is $V$ isomorf met de $K$-vectorruimte $K^n$.
\end{gevolg}
\begin{proof}
    De K-vectorruimte \(K^n\) heeft dimensie \(n\). Het bewijs volgt nu uit gevolg 3.3.5.
\end{proof}

\begin{stelling}[3.3.7 (Dimensiestelling voor lineaire afbeeldingen)]
    Zij $V$ een eindig-dimensionale $K$-vectorruimte, $W$ een willekeurige $K$-vectorruimte,
    en $f \colon V \to W$ een lineaire afbeelding. Dan is
    \[ \dim V = \dim \ker f + \dim \operatorname{im} f. \]
\end{stelling}
\begin{proof}
    Stel basis \(\mathcal{B} = (b_1,\dots, b_m,\dots,b_n)\) voor \(V\) en  basis \(\mathcal{C} = (c_1,\dots,c_m)\) voor \(W\).
    Dan is \(\dim(V) = |\mathcal{B}| = n\) en \(\dim(W) = |\mathcal{C}| = m\), met \(m \leq n\).

    \(\forall v \in V: v = \sum_{i}^{n} \lambda_i b_i \implies f(v) = \sum_{i}^{n} \lambda_i f(b_i) = w \in W \).

    \(\ker f 
    = \{v \in V | f(v) = 0\}
    = \{\lambda_i \in K | \sum_{i}^{n} \lambda_i f(b_i) = 0\}
    \).

    We gebruiken lemma 3.2.6, dat \(f\) volledig bepaald is door de beelden \(f(b)\). 

    Dit impliceerd dat als er \(0 \neq v \in V: \sum_{i}^{n} \lambda_i f(b_i) = 0 \), dat dan \(\lambda_i f(b_i) =  \sum_{j\neq i}^{n-1}\lambda_j f(b_j)\) een lineaire combinatie is van de andere beelden \(f(b_j)\) van de basiselementen.
    Nu gebruiken we gevolg 2.4.8 (ii) \(n-m\) keer tot we een basis krijgen met \(m\) elementen.

    Bijgevolg is \(\dim\ker f = n-m\).

    \(\operatorname{im} f
     = \{f(v)|v \in V\}
     = \{\sum_{i}^{n} \lambda_i f(b_i)|\lambda_i \in K\}
     = \operatorname{span}(\{f(b)|b\in \mathcal{B}\})
     = \operatorname{span}(\{w|w\in W\}).
     \)
     Zodat \(\dim\operatorname{im} f = \dim W = |\mathcal{C}| = m\).

     Uiteindelijk is \(\dim V = |\mathcal{B}| = n \) en \(\dim\ker f + \dim\operatorname{im} f = n-m+m = n \).

     En volgt:
     \[ \dim V = \dim \ker f + \dim \operatorname{im} f.\]

    \textbf{Notebook LM geeft heel prachtig mijn bewijs neergehaald door te zegggen dat \(W\) niet persee eindig-dimensionaal is.}
\end{proof}

\begin{gevolg}[3.3.8 (Alternatief-stelling)]
    Zij $V, W$ $K$-vectorruimten met $\dim V = \dim W < \infty$, en zij $f \colon V \to W$ een lineaire afbeelding. De volgende eigenschappen zijn equivalent:
    \begin{enumerate}
        \item[(a)] $f$ is injectief;
        \item[(b)] $f$ is surjectief;
        \item[(c)] $f$ is bijectief.
    \end{enumerate}
\end{gevolg}
\begin{proof}
    Als \(f\) bijectief is, dan volgt per def dat \(f\) oook surjectief en injectief is.
    Nu hoeven we enkel nog \(f\) injectief \(\iff f\) surjectief aan te tonen.

    \(
        f \text{ bijectief} 
        \iff \ker f = \{0\}
        \iff \dim\ker f = 0
        \iff \dim V = \dim\operatorname{im}f
        \iff \dim W = \dim\operatorname{im}f
        \iff \operatorname{im}f=W
        \iff f \text{ surjectief}
    \).
\end{proof}
\subsection{Optelling en scalaire vermenigvuldiging van lineaire afbeeldingen}

\begin{stelling}[3.4.5]
    Zij $K$ een veld, en $V, W$ twee eindig-dimensionale $K$-vectorruimten.
    Dan is $\dim \operatorname{Hom}(V, W) = \dim V \cdot \dim W$.
\end{stelling}
\begin{proof}
    $$ \operatorname{Hom}_K(V,W) := \{f \colon V \to W \mid f \text{ is een lineaire afbeelding}\} $$
    
    Zij $n = \dim V = |\mathcal{B}|$ met basis $\mathcal{B} = (b_1, \dots, b_n)$ en $m = \dim W = |\mathcal{C}|$ met basis $\mathcal{C} = (c_1, \dots, c_m)$.

    Een lineaire afbeelding $f$ ligt volledig vast door de beelden van de basisvectoren van het domein, namelijk $(f(b_1), \dots, f(b_n))$
    (lemma 3.4.5).

    Het beeld van elke basisvector $b_i$ kan geschreven worden als een lineaire combinatie van de basisvectoren in $W$:
    $$ f(b_i) = \sum_{j=1}^{m} \mu_{ij} c_j $$

    Neem nu een willekeurige vector $v \in V$, die we schrijven als $v = \sum_{i=1}^{n} \lambda_i b_i$. Wegens de lineariteit van $f$ geldt:
    $$ f(v) = \sum_{i=1}^{n} \lambda_i f(b_i) $$
    
    Als we de uitdrukking voor $f(b_i)$ hierin substitueren, krijgen we:
    \[
        f(v) = \sum_{i=1}^{n} \lambda_i \left( \sum_{j=1}^{m} \mu_{ij} c_j \right) 
        \implies f(v) = \sum_{i=1}^{n} \sum_{j=1}^{m} \lambda_i \mu_{ij} c_j
    \]
    Omdat elke lineaire afbeelding $f$ op deze manier volledig en uniek bepaald wordt door de $n \cdot m$ scalairen $\mu_{ij}$, volgt hieruit dat:
    $$ \dim(\operatorname{Hom}(V,W)) = n \cdot m = \dim V \cdot \dim W $$.
\end{proof}
\subsection{De samenstelling van lineaire afbeeldingen}

\begin{lemma}[3.5.2]
    Zij $V, W, U$ drie $K$-vectorruimten, zij $g \in \operatorname{Hom}(V, W)$ en $f \in \operatorname{Hom}(W, U)$. Dan is $f \circ g \in \operatorname{Hom}(V, U)$.
\end{lemma}
\begin{proof}
    We weten dat voor \(v \in V, \, w \in W ,\, u \in U\):

    \(g(v) = w \text{ en }f(w) = u\).

    \((f \circ g)(v) = f(g(v)) = f(w) = u\).

    Als \(\lambda,\mu \in K\) en \(v,v' \in V\):

    \( 
        (f \circ g)(\lambda v + \mu v') 
        = f(g(\lambda v + \mu v'))  
        = f(\lambda g(v) + \mu g(v'))
        = \lambda f(g(v)) + \mu f(g(v'))  
        = \lambda (f \circ g)(v) + \mu (f \circ g)(v') 
    \).

    Hieruit volgt dat \(f \circ g\) lineair is, dus \(f \circ g \in \operatorname{Hom}(V, U)\).
\end{proof}
\subsection{Inverteerbare lineaire operatoren}

\begin{lemma}[3.6.2]
    Zij $f \colon A \to B$ een afbeelding. Dan zijn de volgende uitspraken equivalent:
    \begin{enumerate}
        \item[(a)] Er bestaat een afbeelding $g \colon B \to A$ zodat $f \circ g = \mathbf{1}_B$ en $g \circ f = \mathbf{1}_A$.
        \item[(b)] $f$ is bijectief.
    \end{enumerate}
    Indien deze uitspraken voldaan zijn, dan is $g = f^{-1}$.
\end{lemma}
\begin{proof}
      \([\Rightarrow]\)  Stel er bestaat een afbeelding $g \colon B \to A$ zodat $f \circ g = \mathbf{1}_B$ en $g \circ f = \mathbf{1}_A$.
        \(\forall a \in A: g(f(a)) = \mathbf{1}_A.\) en \(\forall b \in B: f(g(b)) = \mathbf{1}_B.\)

        \(f\) is injectief, want \(f(a)=f(a') \implies g(f(a))=g(f(a')) \implies \mathbf{1}_A(a) = \mathbf{1}_A(a').  \)

        \(f\) is surjectief, want \(b \in B: b = \mathbf{1}_B = f(g(b))\),
        dus \(b\) is het beeld onder \(f\) van het element \(g(b) \in A\).

    \([\Leftarrow]\) Veronderstel omgekeerd, \(f\) bijectief.

    Beschouw \(f^{-1}: B \to A\) met \(f(a) = b\), dan is \(f(f^{-1}(b))=b\forall b \in B\),
    en \(f^{-1}(f(a))=a\forall a \in A\).
\end{proof}